\section{Some problem we met in experiment}
\begin{itemize}
    \item[Q1.] We use x-y mode to observe the dips in experiments 2 and 3, but the data we've saved only got one data point, it doesn't show the whole diagram but a single point.
    \item[A1.] We redo the experiments and use rolling mode rather than x-y mode, then we have whole diagram to analyze.

    \item[Q2.] In the first experiment, we have to observe light intensity under each cell's temperature. But the temperature is not so precise and have to be correction, so we can't use the data when temperature just reach the target temperature, or it will have significantly error.
    \item[A2.] Our solution is to set $5^\circ C$ as a unit to increase or decrease the temperature. We don't just measure the "Heating process" but also measure the "Cooling process".
    So, for example, if we want to record the light intensity under $62^\circ C$, we first raise the temperature from $60^\circ C$ to $65^\circ C$. During the heating process, we take the data of $62^\circ C$, then when the temperature got to $65^\circ C$, we turn to do the cooling process and record the data of $62^\circ C$. Finally, we average these two data to get a relatively precise data.
\end{itemize}

\subsection*{Selected Berkeley pre-lab questions}
We answer three of the Berkeley OPT pre-lab questions~\cite{OPTOpticalPumping2024}
that are most directly supported by our own data.
\begin{itemize}
    \item[Q15.] \emph{At $i = \SI{1}{\ampere}$ through the Helmholtz coils, what Zeeman resonance frequencies are expected for $^{85}$Rb and $^{87}$Rb?}
    \item[A15.] Using the manual's Eq.~(8) with $N = 135$, $a = \SI{27.5}{\centi\meter}$ gives $B(\SI{1}{\ampere}) = \SI{4.42}{\gauss}$; combined with the low-field Breit--Rabi relation $\nu/B = 2.799/(2I+1)$ MHz/G this predicts $\nu_{87} \approx \SI{3.09}{\mega\hertz}$ and $\nu_{85} \approx \SI{2.06}{\mega\hertz}$. Our apparatus uses a different coil pair, so we tested the prediction through the measured main-field calibration (Eq.~\ref{eq:main_cal}): at $I_{\mathrm{main}} = \SI{1}{\ampere}$ the field is $B_{\mathrm{main}} = 8.579(24)\,\si{\gauss}$, giving $\nu_{87}^{\mathrm{pred}} = 6.004(17)\,\si{\mega\hertz}$ and $\nu_{85}^{\mathrm{pred}} = 4.003(11)\,\si{\mega\hertz}$. The $\nu(I_{\mathrm{sweep}})$ fits of Eqs.~\ref{eq:fit_steep}--\ref{eq:fit_shallow} instead slice the field budget at the sweep-only coil ($k_{\mathrm{sweep}} = 0.6397(39)\,\si{\gauss\per\ampere}$), reproducing the same $g_{F,87}/g_{F,85} = 3/2$ ratio with the slope ratio $0.4634/0.2981 = 1.554(22)$ that uniquely pins the two nuclear spins.
    \item[Q16.] \emph{What is the ambient magnetic field in the laboratory?}
    \item[A16.] The Berkeley manual expects the geomagnetic field in Berkeley ($\sim 38^\circ$\,N) to be $\sim\SI{0.48}{\gauss}$ with an inclination of $\sim 60^\circ$. Our apparatus sits in Taipei ($\sim 25^\circ$\,N), where the reference value is $\sim\SI{0.44}{\gauss}$ with inclination $\sim 35^\circ$, giving a horizontal component of $\sim\SI{0.36}{\gauss}$. Two independent measurements fix the component along our optical axis: (i)~the intercept of the combined sweep-field calibration, $B_0 = -0.2263(24)\,\si{\gauss}$ in Eq.~\ref{eq:sweep_cal}; and (ii)~the zero-field resonance current $I_{\mathrm{sweep}}^{(0)} = 0.3567(58)\,\si{\ampere}$ (Rb-85) / $0.3575(80)\,\si{\ampere}$ (Rb-87) shared by both isotopes, which implies $|B_0| = k_{\mathrm{sweep}}\,I_{\mathrm{sweep}}^{(0)} = 0.2282(39)\,\si{\gauss}$. The two agree at the permille level. The residual is smaller than the full horizontal geomagnetic component because the light-proof box is not aligned with magnetic north, and because of partial shielding by the laboratory steelwork; the difference between $\SI{0.23}{\gauss}$ and the $\SI{0.36}{\gauss}$ horizontal reference is the projection-plus-shielding factor.
    \item[Q17.] \emph{What effects would a misalignment between the Helmholtz and ambient fields have on the measurements?}
    \item[A17.] If the ambient field has a transverse component $B_\perp$ that the Helmholtz coil cannot null, the total field along the optical axis is $B_z(I) = k_{\mathrm{sweep}}I + B_{\parallel,0}$ while the field magnitude that the atoms see is $\sqrt{B_z^2 + B_\perp^2}$. Two signatures of this appear in our 4b/4c data. (i)~\emph{Asymmetric intercepts under current reversal.} If $B_\perp \neq 0$, the $\nu(I)$ curves for $+I$ and $-I$ are not related by a simple sign flip: the curves become hyperbolic rather than V-shaped near $I = 0$, with a minimum $\nu_{\min} = (g_F \mu_B / h)\,|B_\perp|$ that is bounded away from zero. Our fitted intercepts, $I_{\mathrm{sweep}}^{(0)} = 0.3575(80)$ and $0.3567(58)$\,A for the two isotopes (Section~\ref{sec:sweep_cal}), agree to one part in $\sim 400$, and the residual field $B_0 = -0.2263(24)\,\si{\gauss}$ is identical between positive- and negative-polarity fits within the quoted error. Both observations bound $|B_\perp|$ to no more than a few mG, well below the \SI{3}{\milli\gauss} digitisation noise of the CH3 axis. (ii)~\emph{The zero-field resonance remains visible.} Even when the axial component is nulled, the surviving transverse residual prevents the $m_F$ sublevels from becoming exactly degenerate; this is precisely why the RF-off trace of Fig.~\ref{fig:4b_raw_traces} still shows a sharp absorption feature rather than the single saturated floor that perfect cancellation would produce. The width of that feature is an upper bound on $|B_\perp|$ and is consistent with the $\sim\SI{0.1}{\mega\hertz}$ Zeeman broadening expected from a few-mG residual.
\end{itemize}